% introduction.tex
Stochastic attention (SA) generates protein sequences from small family alignments by treating
the modern Hopfield energy as an exact score function for Langevin
dynamics~\cite{varnerSAProtein2026}. Given a memory matrix of $K$ PCA-encoded family members, SA
samples from the corresponding Boltzmann distribution without training, producing sequences that
preserve family-level amino acid composition, per-position conservation, and predicted
three-dimensional fold. The method requires only two matrix-vector products per iteration
($\mathcal{O}(dK)$ cost), runs on a laptop in seconds, and is designed for the scarce-data
regime (families with tens to hundreds of members) where learned generative models lack
sufficient training signal. A key limitation of the current framework, however, is that it
treats all stored patterns equally. Every family member contributes the same weight to the
energy landscape, so the Langevin chain explores the full family distribution without preference
for any functional subset. In practice, experimentalists often want something more specific: not
just ``a plausible Kunitz domain,'' but a Kunitz domain that \emph{inhibits trypsin}; not just
``a plausible SH3 domain,'' but one that binds a particular proline-rich motif. Family
alignments define a shared fold (all 99 Kunitz sequences share the same disulfide framework
and double-loop architecture), but the members differ in which target they bind, and that
specificity is encoded in a small number of positions (e.g., the P1 contact residue) rather
than in the global fold. The alignment captures structure; the practitioner wants to condition
on function. The question is: \emph{how can we condition
the generative process on a functional subset without retraining?}

A seemingly simple answer is to \emph{curate} the memory matrix, restricting it to sequences that
a practitioner has identified (through experimental screens, literature annotation, or functional
assays) as belonging to the functional subset of interest. SA then generates new sequences in their
image. We show that this strategy works: even three designated input sequences
suffice for complete phenotype transfer in the Kunitz domain. But hard curation discards all
remaining family members, losing the fold-level constraints they provide. A more principled
approach would retain the full family in the memory while continuously upweighting the designated
subset. The key observation we exploit is that pattern multiplicity in the Hopfield energy provides
exactly this capability. Assigning a multiplicity weight $r_k > 0$ to each stored pattern $\vm_k$
defines a tilted energy landscape whose score function remains analytic. The only modification to
the SA update is adding $\log r_k$ to each logit before the softmax. A single parameter, the
multiplicity ratio $\rho = r_{\mathrm{designated}}/r_{\mathrm{background}}$, parameterizes a
continuous family of Boltzmann distributions that interpolate between unconditioned generation
($\rho=1$, standard SA) and hard subset curation ($\rho \to \infty$). The memory matrix stays
compact at $d \times K$; no columns are duplicated, no parameters are trained, no architectural
changes are needed. The method is agnostic to what the designated subset represents (binding,
thermostability, organism of origin); it simply biases the sampler to spend more time near the
patterns the user points to.

In this study, we developed the theory of multiplicity-weighted SA and tested it on five Pfam
protein families (Kunitz, SH3, WW, Homeobox, and Forkhead domains) spanning a range of family
sizes ($K = 55$--$420$) and functional-subset geometries. We showed that the multiplicity
weights enter the sampler as a simple logit bias, preserving the analytic score function and
$\mathcal{O}(dK)$ cost of the original method. The energy-level conditioning is exact: the
softmax attention on designated patterns matches the target fraction to within $0.3\%$ at all
multiplicity ratios tested. However, we found that the decoded sequence phenotype can fall short
of this target, because the PCA encoding that compresses sequences into the sampler's state space
does not always preserve the residue-level variation that defines the functional split. We termed
this discrepancy the calibration gap and showed that it is predicted by a simple geometric
quantity, the Fisher separation index $S$, which measures how well PCA separates designated from
background sequences. Across five families spanning a fourfold range of $S$, the gap decreased
monotonically with increasing separation ($R^{2} = 0.80$). This provides a practitioner's
decision rule: compute $S$ before generating any sequences; if $S$ is high, multiplicity
weighting transfers the phenotype directly; if $S$ is low, use hard curation instead.
